\section{Why the Lambek calculus matters}
The Lambek calculus was introduced in \citet{Lambek1958} as a formal calculus of syntactic
types. Its original purpose was already recognisably linguistic: to give a principled method for
distinguishing well-formed from ill-formed strings and to model the way in which syntactic
categories combine in natural language. Lambek's follow-up work in \citet{Lambek1961}
further developed the type-theoretic viewpoint and helped establish categorial grammar as a
serious mathematical approach to syntax.

The importance of the Lambek calculus lies in the fact that it sits at an unusually fruitful
intersection of logic, algebra, and linguistics. On the logical side, it is one of the prototypical
substructural logics: order matters, and structural principles such as exchange, weakening, and
contraction are not built in by default. On the linguistic side, this resource sensitivity is not an
artificial restriction but an advantage, because word order and local composition are central to
syntactic analysis. For this reason the Lambek calculus became a foundational system in the
study of categorial and type-logical grammars \citep{vanBenthem1988,Moortgat1995Multimodal,morrill1996generalising}.

A second reason for its importance is that the calculus offers an explicit bridge between syntax
and semantics. Derivations in categorial calculi naturally correspond to compositional analyses of
meaning, and this connection has been one of the main reasons for the enduring influence of
Lambek-style systems in formal semantics and mathematical linguistics \citep{vanBenthem1988,Steedman2000}.

